\documentclass[a4paper,12pt]{report}
\usepackage[utf8x]{inputenc}
\usepackage[francais]{babel}
\usepackage[top=2cm, bottom=2cm, left=2cm, right=2cm]{geometry}
\usepackage{graphicx}
\usepackage{float}
\usepackage{listings} %Pour le code
\usepackage{ifthen}
\usepackage{color}
\usepackage{textcomp}


\title{Rapport du projet informatique}
\author{Matthieu Denoux - Groupe 1}


\begin{document}
\maketitle
\tableofcontents
\newpage

\part{Analyse du travail effectué}
%%%%%%%%%%%%%%%%%%%%%%%%%%%%%%%%%%
%%                              %%
%%   ------------------------   %%
%%                              %%
%%%%%%%%%%%%%%%%%%%%%%%%%%%%%%%%%%
\chapter{Présentation du sujet}

%%%%%%%%%%%%%%%%%%%%%%%%%%%%%%%%%%
%%   ------------------------   %%
%%%%%%%%%%%%%%%%%%%%%%%%%%%%%%%%%%
\section{Principe général}
\paragraph{}
Le sujet choisi a pour intitulé \textbf{Babyfoot en réseau}. Il s'agit de concevoir un système complet de jeu en réseau.
Le système serait donc séparé en deux parties, un serveur et un client. Chaque joueur pourrait donc se connecter à une partie
n'ayant pas encore commencé et une fois le nombre de joueurs réunis (\emph{2, 3 ou 4}), la partie serait lancée. Il faut donc réaliser à la fois 
le système réseau, l'interface graphique et imaginer un gameplay qui rende le jeu agréable.

%%%%%%%%%%%%%%%%%%%%%%%%%%%%%%%%%%
%%   ------------------------   %%
%%%%%%%%%%%%%%%%%%%%%%%%%%%%%%%%%%
\section{L'interface graphique}
\paragraph{}
Au niveau de l'interface graphique, il faudra réaliser plusieurs fenêtres successives de menus pour parvenir jusqu'au jeu lui-même. Ces fenêtres seront réalisées avec la bibliothèque inclue dans le package standard \textbf{Swing}. J'ai aussi choisi d'utiliser des JFrame et donc de programmer une application à part entière et non une appliquette.
\paragraph{Menu principal}
Un premier menu, dit \emph{menu principal}, permettra de commencer une nouvelle partie en ligne, de rejoindre une partie
déjà en cours, de modifier les quelques options disponibles ou bien de quitter le jeu. Un header unifie les différentes fenêtres, il donnera en plus une certaine identité graphique au jeu avec l'image affichée.
    
    \begin{figure}[H]
	\begin{center}
	\includegraphics[width=17cm]{pictures/header.png} 
	\end{center}
	\caption{Header du programme, très peu de recherche esthétique, simplement une image de base.}
    \end{figure}
\paragraph{Commencer une partie}
Si l'on commence une nouvelle partie en ligne, on se retrouve dans une \og salle de configuration \fg. Là, il est possible de configurer la partie que l'on souhaite lancer, l'ouvrir à d'autres joueurs, ce qui nous fait arriver dans une \og salle d'attente \fg puis d'attendre que d'autres joueurs rejoignent la partie. Une fois que les équipes sont complètes, on peut lancer le jeu.
\paragraph{Rejoindre une partie}
On peut aussi sélectionner une partie dans la liste des parties déjà commencées, dans la limite des places disponibles. On rejoint alors une \og salle d'attente \fg semblable à celle décrit dans la section précédente où l'on attend que la partie soit complète avant que le meneur, celui qui a créé la partie, ne lance le jeu. C'est le seul qui dispose de cette possibilité.
\paragraph{Déroulement d'une partie}
Il y aurait deux gameplays différents : 
\begin{itemize}
    \item l'un entièrement manuel verrait le joueur maître de ses possibilités. Ainsi, les touches \og A \fg, \og Z \fg, \og E \fg, \og R \fg permettraient de sélectionner la canne que l'on souhaiterait manier. Le déplacement de la souris latéralement permettraient quant à elle de déplacer les cannes tandis que la barre espace commanderait le tir. Je ne pense pas réaliser plusieurs puissances de tir possibles mais cela pourrait consister en une amélioration de la richesse du gameplay.
    \item l'autre, principalement rencontrée dans les autres jeux de babyfoot trouvés en ligne, serait plus automatisée : les cannes se déplaceraient toutes ensemble de manière synchronisée avec les déplacements de la souris. Le tir pourrait être ou non automatisé selon que j'aurais le temps pour configurer cette fonctionnalité supplémentaire.
\end{itemize}
\paragraph{Options}
Un petit menu sera consacré aux options, notamment le type de jeu que l'on souhaite utiliser (manuel ou automatisé, auquel cas jusqu'à quel point).


%%%%%%%%%%%%%%%%%%%%%%%%%%%%%%%%%%
%%                              %%
%%   ------------------------   %%
%%                              %%
%%%%%%%%%%%%%%%%%%%%%%%%%%%%%%%%%%
\chapter{Analyse de la solution envisagée}


%%%%%%%%%%%%%%%%%%%%%%%%%%%%%%%%%%
%%   ------------------------   %%
%%%%%%%%%%%%%%%%%%%%%%%%%%%%%%%%%%
\section{Découpage en modules}
\paragraph{}
Je prévois de découper le code en deux grandes parties :
\begin{itemize}
    \item Partie client
    \item Partie serveur
\end{itemize}
Les deux parties (client/serveur) seront elles-mêmes séparées en plusieurs modules.
\begin{itemize}
    \item L'interface graphique (rangée dans /gui) qui sera absente du serveur
    \item La partie réseau de l'application (rangée dans /network), responsable de l'émission et de la réception des requêtes.
    \item Le c\oe{}ur algorithmique de l'application (rangée dans /core) contenant les classes gérant les matchs, les joueurs, etc.
\end{itemize}
Chaque partie et chaque module comporteront donc plusieurs classes qui se chargeront chacune d'une des tâches du module parent.


%%%%%%%%%%%%%%%%%%%%%%%%%%%%%%%%%%
%%   ------------------------   %%
%%%%%%%%%%%%%%%%%%%%%%%%%%%%%%%%%%
\section{L'interface graphique}
\paragraph{MainFrame.java}
Gère la fenêtre qui englobe tout le reste. On utilisera en fait des JPanel pour modifier le contenu de cette fenêtre. J'empêcherai dans un premier temps de modifier la taille de la fenêtre pour éviter d'avoir des problèmes de dessin du terrain de babyfoot à gérer. Cette classe contient le \textbf{main} du programme. Elle permet d'accéder aux différents éléments : le client récepteur, les informations sur le joueur lui-même, etc.
\paragraph{BPanel.java}
Classe abstraite présentant certaines actions mécaniques quant aux caractéristiques des JPanel utilisés (header, background color, taille, layoutmanager, etc.). Elle contiendra notamment une référence vers l'instance de la fenêtre de type MainFrame La plupart des panels ci-dessous en héritent et peuvent donc accéder à cette classe.
\paragraph{ConnexionPanel.java}
Gère le premier écran affiché à l'ouverture du jeu. On y choisit le pseudonyme que l'on souhaite utiliser dans le reste du jeu et qui s'affichera ensuite sur les écrans lors du contact avec d'autres joueurs. Un message d'erreur est affiché en rouge lorsque le pseudo est déjà utilisé.
\paragraph{MenuPanel.java}
Gère le menu principal affiché après connexion. Contient les boutons qui mèneront vers les principales actions citées plus haut (nouvelle partie, rejoindre, options, quitter).
\paragraph{NewPanel.java}
Lorsque l'on lance une nouvelle partie, on obtient cet écran qui contiendra les principales options nécessaires pour configurer une partie puis l'ouvrir à des joueurs extérieurs. Une fois cette partie ouverte, on aboutit à une \emph{salle d'attente} gérée par WaitingRoomPanel.
\paragraph{ServersPanel.java}
Gère la liste des parties actuellement en recherche de joueurs lorsque l'on cherche à rejoindre une partie. On sélectionne une partie dans cette liste puis on aboutit à une \emph{salle d'attente}.
\paragraph{WaitingRoomPanel.java}
Une fois la partie configurée et ouverte, on aboutit dans cette salle d'attente qui attendra que certaines conditions soient réunies pour permettre au jeu d'être commencé. On peut \emph{aussi} y accéder depuis la partie \emph{rejoindre une partie}.
\paragraph{GamePanel.java}
La partie à proprement parler. Donc le conteneur du dessin du terrain qui gère les éléments extérieurs, les événements et toute autre interaction avec le reste du code. Inclut la zone de dessin décrite ci-dessous.
\paragraph{GameZone.java}
Zone de dessin chargée de représenter le terrain, les joueurs, etc. Fichier assez conséquent puisqu'il contient toute la gestion de l'affichage et des données liées aux déplacements. Toutes les informations, les calculs sont effectués par le serveur et envoyées.
\paragraph{ChatPanel.java}
Une partie un peu à part qui est chargée de gérer l'affichage du chat. Celui-ci sera présent à de nombreux endroits : la liste des parties disponibles, la salle d'attente, la création d'une partie et durant le jeu lui-même. Or, tous ces affichages sont centralisés dans cet unique fichier qui centralise ainsi le traitement et l'affichage et permet de réutiliser à plusieurs endroits ces données.


%%%%%%%%%%%%%%%%%%%%%%%%%%%%%%%%%%
%%   ------------------------   %%
%%%%%%%%%%%%%%%%%%%%%%%%%%%%%%%%%%
\section{Le réseau}
\subsection{Le principe général}
\paragraph{}
Des messages sont échangés entre le client et le serveur et ce de façon répétée. Il faut donc utiliser des Thread pour gérer l'émission et la réception de données depuis le serveur et un Thread pour gérer la réception côté client. Le serveur se chargera ensuite d'enregistrer les données statiques devant être conservées.
\subsection{La partie client}
\paragraph{ClientBabyfoot.java}
Il s'agit du main du client qui initialise les fenêtres et les différentes classes utilisées.
\begin{quote}
	private static Player player;
    
	private static Chat chat;
    
	private static Client client;
\end{quote}
\paragraph{Client.java}
Cette classe gère la connexion au serveur : il y a en tout trois connexions (une pour le tchat, une pour les joueurs et une pour les données des matchs) afin de pouvoir envoyer plusieurs requêtes simultanément du même client. Toutes les requêtes envoyées sont sous la forme de chaîne de caractères avec l'utilisation d'un caractère spécial, qui est contenu dans une constante de la classe Utils du Core comme séparateur.
\begin{quote}
	private static Socket socketChat;
	
	private static Socket socketPlayer;
	
	private static Socket socketMatch;
	
	private static Socket socketGame;
	
	private ChatClient cc;
	
	private PlayerClient pc;
	
	private MatchClient mc;
	
	private GameClient gc;
	
	private Thread tChat;
	
	private Thread tPlayer;
	
	private Thread tMatch;
	
	private Thread tGame;
\end{quote}
\paragraph{PlayerClient.java}
Gère les différentes actions possibles par et sur les joueurs et les requêtes envoyées au serveur pour pouvoir satisfaire le système (connexion d'un joueur, déconnecter un joueur, ajouter un match, etc.). Elle contient aussi un Thread qui écoute l'entrée de la socket.
\paragraph{MatchClient.java}
Procède de la même façon que PlayerClient pour les matchs, envoie les requêtes au serveur et récupère les données sur la position de la balle, des barres des différents joueurs, etc.
\paragraph{ChatClient.java}
S'occupe de gérer la partie client du tchat : une fois instanciée par le ChatPanel, elle permet d'envoyer des données au serveur et de récupérer la liste des messages disponibles dans le salon. Elle permet aussi de changer de salon, d'afficher la liste des salons et celle des joueurs connectés sur le jeu.
\subsection{La partie serveur}
\paragraph{AbstractServer.java}
Classe abstraite qui gère quelques méthodes de base que doivent toutes présenter les classes gérant les serveurs.
\paragraph{ServerBabyfoot.java}
Gère les différentes requêtes et les redirige vers les autres entités du serveur (chat, jeu, etc.). Contient un système de Thread pour pouvoir gérer simultanément ces différentes actions, un Thread général qui gère les requêtes. C'est aussi le main qui lance le serveur.
\begin{quote}
    ServerSocket socketserver = null;
    
	Socket socket = null;
    
	BufferedReader in;
    
	PrintWriter out;
    
	String login;
    
	public static ChatServer tchat;
    
	public static PlayerServer tplayer;
    
	public static MatchServer tmatch;
\end{quote}
\paragraph{PlayerServer.java}
Gère la connexion au serveur d'un des joueurs. L'enregistre dans une liste de joueurs connectés. Peut aussi renvoyer la liste des joueurs connectés actuellement et leur état.
\paragraph{MatchServer.java}
Gère tous les échanges de données sur les matchs. S'occupe ensuite de fournir les informations nécessaires au client sur l'état actuel de la position de la balle par exemple. Nécessitera peut-être, pour des questions de réactivité, d'être multi-threadé.
\paragraph{ChatServer.java}
S'occupe de gérer la partie serveur du tchat : écoute un port, récupère les messages envoyés par des joueurs ainsi que le salon où ils se trouvaient à ce moment là. Les enregistre dans un fichier ou bien dans une base de données (à voir).


%%%%%%%%%%%%%%%%%%%%%%%%%%%%%%%%%%
%%   ------------------------   %%
%%%%%%%%%%%%%%%%%%%%%%%%%%%%%%%%%%
\section{Le c\oe{}ur algorithmique}
\paragraph{Player.java}
Le bloc de base du c\oe{}ur algorithmique est le joueur. Il est instancié par le serveur et a plusieurs attributs : son login, son état actuel (s'il est en train de jouer ou non). Si le joueur est en match, la classe contient aussi les barres qu'il a le droit de déplacer.
\begin{quote}
	private Match match;
    
	private String login;
\end{quote}
\paragraph{Utils.java}
Est une classe abstraite qui gère les fonctions utilitaires comme la mise-en-forme de la date, la mise-en-forme des requêtes, les fonctions de hash utilisées pour vérifier la validité des requêtes, etc. Commune au client et au serveur.
\paragraph{Match.java}
Contient les données principales pour une partie, notamment l'avancement de la partie, l'état des scores, les joueurs y participant, etc. Sera appelée par le serveur. Cette classe contient aussi toutes les informations sur les données factuelles d'une partie, à savoir la position de la balle, la position des barres, etc. Elle se charge d'effectuer l'ensemble des calculs de déplacements de la balle, comptabilise les buts, met en place les pauses.
\paragraph{Collisions.java}
Gère tout ce qui touche aux collisions, avec les joueurs ou les bords. C'est là que se situe toute la difficulté de calcul.
\begin{quote}
    private int leftScore;
    
	private int rightScore;
	
	private Types type;
	
	private States state;
	
	private Player player1;
    
	private Player player2;
    
	private Player player3;
    
	private Player player4;
\end{quote}
\paragraph{Database.java}
Aucune base de données n'a finalement été utilisée. Le serveur stocke donc ses données dans la RAM puis les supprime lors de sa fermeture. Dans un souci de pérennité et de performances, il serait intéressant de les stocker dans des fichiers textes que l'on puisse ainsi conserver des archives.
%%%%%%%%%%%%%%%%%%%%%%%%%%%%%%%%%%
%%                              %%
%%   ------------------------   %%
%%                              %%
%%%%%%%%%%%%%%%%%%%%%%%%%%%%%%%%%%
\chapter{Échéancier}
\paragraph{Ce qui était prévu}
\begin{enumerate}
    \item \textbf{Fin Décembre} : rédiger l'avant-projet.
    \item \textbf{Début/Mi Janvier} : obtenir la validation et les annotations sur la structure choisie.
    \item \textbf{Fin Janvier} : réaliser la partie graphique du programme et avoir regardé les grandes lignes du développement serveur / bases de données. Avoir mis au point les éléments de base du gameplay (interaction joueur/machine).
    \item \textbf{Fin Février} : développer le serveur et la gestion des différents types de requêtes. Mettre au point le tchat et la gestion des joueurs avec la base de données.
    \item \textbf{Mi Mars} : dresser les liens entre serveur et jeu. Tester.
\end{enumerate}
\paragraph{Ce qui a été fait}
\begin{enumerate}
    \item \textbf{Fin Décembre} : rédiger l'avant-projet.
    \item \textbf{Fin décembre/Début Janvier} : développer la partie réseau et l'architecture de base des requêtes.
    \item \textbf{Fin Janvier} : réalisation de l'interface graphique et des premiers calculs pour la mise-en-place de matchs.
    \item \textbf{Février} : ajout des collisions et corrections des principaux bugs liés à la gestion des joueurs, matchs et chat.
    \item \textbf{Mars} : débuguer le reste du programme et ajouter les fonctionnalités supplémentaires nécessaires à un bon fonctionnement.
\end{enumerate}
\part{Explications des choix techniques}
%%%%%%%%%%%%%%%%%%%%%%%%%%%%%%%%%%
%%   ------------------------   %%
%%%%%%%%%%%%%%%%%%%%%%%%%%%%%%%%%%
\section{Le réseau}
\paragraph{}
Le réseau est au coeur de mon projet et j'ai donc développé plusieurs Thread en parallèle pour pouvoir gérer plus facilement l'envoi de messages par mon application.
\section{La gestion du jeu}
\paragraph{}
Pour calculer la position de la balle au cours du temps, j'ai utilisé un Thread qui s'actualise toutes les 10ms et qui effectue à chaque tour d'horloge
l'actualisation de la position à partir de la vitesse et des tests de collisions. Les tests de collisions sont calculés à part dans une classe appelée Collision.
Je vérifie ainsi la position de la balle par rapport aux bords et détermine si une collision a eu lieu avec une de ces limites. Une fonction est appelée aussi afin de
déterminer si il n'y a pas eu de contact avec un des joueurs du terrain.
\paragraph{}
La principale difficulté algorithmique vient de cet aspect justement. Parvenir à prévoir les différents comportements de la balle par rapport aux
joueurs. Les coordonnées utilisées par ces calculs diffèrent en plus parfois de celles d'affichage et il faut prendre en compte la taille de l'écran ainsi que les imprécisions dues au caractère discret du calcul. J'ai encore parfois des problèmes de collisions qui se produisent trop tôt, c'est-à-dire que le joueur humain a l'impression qu'elle se cogne contre un objet invisible, ou encore des collisions gérées trop tard auquel cas la balle rentre partiellement dans le joueur en plastique.

\lstset
{
	numberstyle=\footnotesize,
	basicstyle=\ttfamily\footnotesize,
	numbers=left,
	stepnumber=1,
	showstringspaces=false,
	showspaces=false,
	breaklines=true,
	frame=single,
	framexleftmargin=1mm,
	framextopmargin=1mm
}

\definecolor{gray}{gray}{0.5}
\definecolor{green}{rgb}{0,0.5,0}

\lstloadlanguages{Java}
\lstset{language=Java}
\lstset{%
    inputencoding=utf8,
	extendedchars=true,
	literate=%
	{é}{{\'{e}}}1
	{è}{{\`{e}}}1
	{ê}{{\^{e}}}1
	{ë}{{\¨{e}}}1
	{û}{{\^{u}}}1
	{ù}{{\`{u}}}1
	{â}{{\^{a}}}1
	{à}{{\`{a}}}1
	{î}{{\^{i}}}1
	{ô}{{\^{o}}}1
	{ç}{{\c{c}}}1
	{Ç}{{\c{C}}}1
	{É}{{\'{E}}}1
	{Ê}{{\^{E}}}1
	{À}{{\`{A}}}1
	{Â}{{\^{A}}}1
	{Î}{{\^{I}}}1
}
\lstset
{
	language=java,
	basicstyle=\ttfamily\footnotesize,
	stringstyle=\color{red},
	alsoletter={1234567890},
	otherkeywords={\ , \}, \{},
	keywordstyle=\color{blue},
	upquote=true,
	morecomment=[s]{"""}{"""},
	commentstyle=\color{gray}\slshape
}

\part{Le code partiel du programme}
	\setcounter{chapter}{0}
	\chapter{Client}
	\section{Core}
        \subsection{ClientBabyfoot.java}
            \lstinputlisting[language=Java, firstline=11, lastline=58]{../babyfootClient/clientCore/ClientBabyfoot.java}
        \subsection{Utils.java}
            \lstinputlisting[language=Java, firstline=9, lastline=139]{../babyfootClient/clientCore/Utils.java}
    \section{Gui}
        \subsection{BPanel.java}
            \lstinputlisting[language=Java, firstline=8, lastline=32]{../babyfootClient/clientGui/BPanel.java}
        \subsection{GameZone.java}
            \lstinputlisting[language=Java, firstline=39, lastline=237]{../babyfootClient/clientGui/GameZone.java}
            \lstinputlisting[language=Java, firstline=650, lastline=688]{../babyfootClient/clientGui/GameZone.java}
    \section{Network}
        \subsection{Client.java}
            \lstinputlisting[language=Java, firstline=11, lastline=57]{../babyfootClient/clientNetwork/Client.java}
        \subsection{MatchClient.java}
            \lstinputlisting[language=Java, firstline=15, lastline=340]{../babyfootClient/clientNetwork/MatchClient.java}
    \section{Classe de test}
        \lstinputlisting[language=Java, firstline=12, lastline=53]{../babyfootClient/clientTest/Tester.java}

	\chapter{Serveur}
	\section{Core}
        \subsection{Fonction de tir dans Match.java}
            \lstinputlisting[language=Java, firstline=347, lastline=450]{../babyfootServer/serverCore/Match.java}
        \subsection{Thread de Match.java}
            \lstinputlisting[language=Java, firstline=483, lastline=540]{../babyfootServer/serverCore/Match.java}
        \subsection{Collisions.java}
            \lstinputlisting[language=Java, firstline=16, lastline=214]{../babyfootServer/serverCore/Collisions.java}
    \section{Network}
        \subsection{PlayerServer.java}
            \lstinputlisting[language=Java, firstline=13, lastline=68]{../babyfootServer/serverNetwork/PlayerServer.java}
        \subsection{ServerBabyfoot.java}
            \lstinputlisting[language=Java, firstline=13, lastline=103]{../babyfootServer/serverNetwork/ServerBabyfoot.java}

\end{document}          
